% SGA 3, Expose V, Loop 1 English reconstruction.
% Scope: component-local pp. 39--41, proof of Theorem 10.4.1 through
% Corollary 10.4.3 and the complete bibliography.
% Authority: Polo--Gille Exp5-13oct24.pdf, combined-reader pp. 301--303,
% re-edited headers 289--291 (original running markers 292--294).
% Jacob Reinhold's CC BY 4.0 Markdown is comparison/drafting material only.

\medskip
\phantomsection\label{proof:sga3-V-10.4.1-a}%
\noindent\textit{\hyperref[thm:sga3-V-10.4.1-a]{a)} 的证明.}---
假设 fppf 层 \(G_X/F\) 由一个 \(X\)-概形
\(Y\) 表出. 则由 \SGARefLink{sga3:iv:corollary:6:3:3}{第~IV 篇, 6.3.3}, 态射
\(p:G_X\to Y\) 忠实平坦且局部有限呈示,
而\hyperref[fig:sga3-V-10.4.1-first-cartesian]{下图}中的第二个方块为笛卡尔方块:
\SGAThreeDiagram[0.66\textwidth]
  {figures/exp5/p289_theorem1041_first_cartesian.png}
  {定理~\hyperref[thm:sga3-V-10.4.1-a]{10.4.1(a)}
   的证明中使用的笛卡尔方块.}
  {fig:sga3-V-10.4.1-first-cartesian}
  {第 V 篇, 组件局部第~39 页 / 合订阅读版第~301 页 /
   修订版页眉第~289 页}
\hyperref[fig:sga3-V-10.4.1-first-cartesian]{第一个方块}由沿单位截面
\(e_X:X\to G_X\) 作基变换得到. 由于 \(p\) 忠实平坦且局部
有限呈示, \(F\to X\) 也有同样的性质.

反之, 假设 \(F\) 在 \(X\) 上平坦, 并令
\(X^2=X\times_SX\). 态射 \(d:G_X\to X^2\) 给出纤维积
方块
\SGAThreeDiagram[0.48\textwidth]
  {figures/exp5/p289_theorem1041_fiber_square.png}
  {与 \(d:G_X\to X^2\) 相伴的纤维积方块.}
  {fig:sga3-V-10.4.1-fiber-square}
  {第 V 篇, 组件局部第~39 页 / 合订阅读版第~301 页 /
   修订版页眉第~289 页}
\phantomsection\label{eq:sga3-V-10.4.1-pr1-torsor}%
此时第一投影
\[
  \operatorname{pr}_1:G_X\times_{X^2}G_X\longrightarrow G_X
\]
是 \(G_X\) 上的一个 \(F\times_XG_X\)-挠子, 因而平坦
且有限型, 因为 \(F\) 具有这些性质.%
\footnote{\label{source-note:sga3-V-10.4.1-pr1-torsor}源 PDF 此处却印作态射
\(G_X\times_{X^2}G_X\to X^2\), 并称它为
\(X^2\) 上的 \(F\times_XX^2\)-挠子. 该态射一般并不
满射. 到 \(G_X\) 的\hyperref[eq:sga3-V-10.4.1-pr1-torsor]{第一投影}才是稳定化子挠子,
并提供定理~\hyperref[thm:sga3-V-10.1.1]{10.1.1} 所需的平坦性; 英文版采用
这一类型正确的读法.}
由定理~\hyperref[thm:sga3-V-10.1.1]{10.1.1}, \(d\) 唯一地分解
为
\SGAThreeDiagram[0.47\textwidth]
  {figures/exp5/p289_theorem1041_factorization.png}
  {\(d\) 的唯一分解.}
  {fig:sga3-V-10.4.1-factorization}
  {第 V 篇, 组件局部第~39 页 / 合订阅读版第~301 页 /
   修订版页眉第~289 页}
其中 \(\psi\) 忠实平坦且有限型, \(\tau\) 为
概形的单态射.

因此层态射 \(\psi:G_X\to Y\) 是
\(F\)-不变的, 因而诱导态射
\[
  \bar\psi:G_X/F\longrightarrow Y.
\]
此外, 由于 \(\psi\) 忠实平坦且有限型,
层的单态射 \(\tau\) 经由 \(d\) 的层像,
即 \(\widetilde\Gamma\) 分解. 因此同构
\(G_X/F\simeq\widetilde\Gamma\) 经由单态射
\(Y\to\widetilde\Gamma\) 分解. 由此得出 \(Y\) 表出 \(G_X/F\).

\medskip
\phantomsection\label{proof:sga3-V-10.4.1-b}%
\noindent\textit{\hyperref[thm:sga3-V-10.4.1-b]{b)} 的证明.}---
假设 \(F\) 在 \(X\) 上平坦. 由 \hyperref[thm:sga3-V-10.4.1-a]{(a) 部分}及\hyperref[proof:sga3-V-10.4.1-a]{其证明}, \(G_X/F\)
可表, 且
\[
  p:G_X\longrightarrow G_X/F
\]
忠实平坦且有限呈示. 另一方面,
由假设, 对 \(i=0,1\), 态射 \(d_i:G_X\to X\) 忠实平坦
且有限呈示. 由于
\(d_i=\bar d_i\circ p\), 由 EGA~IV\(_2\), 2.2.13(iii)
和 EGA~IV\(_3\), 11.3.16, 每个 \(\bar d_i\) 都忠实平坦
且有限呈示.

\medskip
\phantomsection\label{thm:sga3-V-10.4.2}%
\noindent\textbf{定理 10.4.2.}---%
\footnote{\label{note:sga3-V-52}编者注: 这里也存在一个
满足\hyperref[thm:sga3-V-10.4.2]{此定理}结论的最大开子集 \(U\subset X\).
此外, 余维一的点 \(x\in X\) 属于
\(U\) 当且仅当
\[
  (G_X/F)\times_X\operatorname{Spec}(\mathcal O_{X,x})
  \longrightarrow
  \operatorname{Spec}(\mathcal O_{X,x})
  \times_S\operatorname{Spec}(\mathcal O_{X,x})
\]
为闭浸入; 见编者\hyperref[note:sga3-V-44]{注~44}.}
\textit{在 \hyperref[sec:sga3-V-10.2]{10.2} 的假设下,
假设 \(F\) 在 \(X\) 上平坦. 则存在一个稠密饱和
开子集 \(U\subset X\), 使得 fppf 商
\(V=G\backslash U\) 为有限型 \(S\)-概形, 且
\(U\to V\) 忠实平坦且有限呈示.}

\medskip
\noindent\textit{证明.}---
定理~\hyperref[thm:sga3-V-10.4.1]{10.4.1} 表明
\(G_X/F\simeq\widetilde\Gamma\) 可表. 因此 fppf 层
\(G\backslash X\) 为下列态射对的商层:
\SGAThreeDiagram[0.38\textwidth]
  {figures/exp5/p290_theorem1042_double_arrow.png}
  {其商为 \(G\backslash X\) 的等价对.}
  {fig:sga3-V-10.4.2-double-arrow}
  {第 V 篇, 组件局部第~40 页 / 合订阅读版第~302 页 /
   修订版页眉第~290 页}
由\hyperref[thm:sga3-V-10.4.1-b]{上述结果}, 对 \(i=0,1\),
\[
  \bar d_i:G_X/F\longrightarrow X
\]
忠实平坦且有限呈示, 而态射
\SGAThreeDiagram[0.53\textwidth]
  {figures/exp5/p290_theorem1042_monomorphism_chain.png}
  {从商 \(G_X/F\) 到 \(X\times_SX\) 的单态射.}
  {fig:sga3-V-10.4.2-monomorphism}
  {第 V 篇, 组件局部第~40 页 / 合订阅读版第~302 页 /
   修订版页眉第~290 页}
为单态射. 换言之, \((\bar d_0,\bar d_1)\) 是一个
等价对. 因此可以应用
定理~\hyperref[thm:sga3-V-8.1]{8.1}. 从而存在一个
稠密饱和开子集 \(U\subset X\), 使得 fppf 商
\(V=G\backslash U\) 是有限型 \(S\)-概形, 且 \(U\to V\)
忠实平坦且有限呈示.

结合泛平坦性定理
(EGA~IV\(_2\), 6.9.3), 得到\hyperref[cor:sga3-V-10.4.3]{下述结果}.

\medskip
\phantomsection\label{cor:sga3-V-10.4.3}%
\noindent\textbf{推论 10.4.3.}---
\textit{在 \hyperref[sec:sga3-V-10.2]{10.2} 的假设下,
假设 \(X\) 约化. 则存在一个稠密饱和开
子集 \(U\subset X\), 使得 fppf 商
\(G\backslash U\) 是有限型 \(S\)-概形, 且
\(U\to G\backslash U\) 忠实平坦且有限
呈示.}

\section*{参考文献\protect\footnotemark}
\phantomsection\label{sec:sga3-V-bibliography}
\addcontentsline{toc}{section}{参考文献}
\footnotetext{\label{note:sga3-V-53}编者注: 本篇引用的补充文献.}

\begin{description}
\item[{[AK80]}]\phantomsection\label{bib:sga3-V-AK80} A.~B. Altman 与 S.~L. Kleiman,
  \emph{Compactifying the Picard Scheme}, Adv. Math. \textbf{35}
  (1980), 50--112.
\item[{[An73]}]\phantomsection\label{bib:sga3-V-An73} S.~Anantharaman,
  \emph{Sch\'emas en groupes, espaces homog\`enes et espaces
  alg\'ebriques sur une base de dimension 1}, M\'em. Soc. Math.
  France \textbf{33} (1973), 5--79.
\item[{[BLR90]}]\phantomsection\label{bib:sga3-V-BLR90} S.~Bosch, W.~L\"utkebohmert 与 M.~Raynaud,
  \emph{N\'eron models}, Springer-Verlag, 1990.
\item[{[CTS79]}]\phantomsection\label{bib:sga3-V-CTS79} J.-L.~Colliot-Th\'el\`ene 与 J.-J.~Sansuc,
  \emph{Fibr\'es quadratiques et composantes connexes r\'eelles},
  Math. Ann. \textbf{244} (1979), 105--134.
\item[{[DG70]}]\phantomsection\label{bib:sga3-V-DG70} M.~Demazure 与 P.~Gabriel,
  \emph{Groupes alg\'ebriques}, Masson \& North-Holland, 1970.
\item[{[DR81]}]\phantomsection\label{bib:sga3-V-DR81} J.~Dixmier 与 M.~Raynaud,
  \emph{Sur le quotient d'une vari\'et\'e alg\'ebrique par un groupe
  alg\'ebrique}, 载于 \emph{Mathematical Analysis and
  Applications}, 第~327--344 页 (L.~Schwartz 65 岁诞辰纪念, L.~Nachbin 编), Adv.
  Math. Suppl. Stud., 第~7A 卷, 1981.
\item[{[Fe03]}]\phantomsection\label{bib:sga3-V-Fe03} D.~Ferrand,
  \emph{Conducteur, descente et pincement}, Bull. Soc. Math. France
  \textbf{131} (2003), 第~4 期, 553--585.
\item[{[Hi62]}]\phantomsection\label{bib:sga3-V-Hi62} H.~Hironaka,
  \emph{An example of a non-K\"ahlerian complex analytic deformation
  of K\"ahlerian complex structures}, Ann. of Math. \textbf{75}
  (1962), 第~1 期, 190--208.
\item[{[KM97]}]\phantomsection\label{bib:sga3-V-KM97} S.~Keel 与 S.~Mori,
  \emph{Quotient by groupoids}, Ann. of Math. \textbf{145} (1997),
  第~1 期, 193--213.
\item[{[Ko97]}]\phantomsection\label{bib:sga3-V-Ko97} J.~Koll\'ar,
  \emph{Quotient spaces modulo algebraic groups}, Ann. of Math.
  \textbf{145} (1997), 第~1 期, 33--79.
\item[{[Ko08]}]\phantomsection\label{bib:sga3-V-Ko08} J.~Koll\'ar,
  \emph{Quotients by finite equivalence relations}, arXiv:0812.3608.
\item[{[Mum65]}]\phantomsection\label{bib:sga3-V-Mum65} D.~Mumford,
  \emph{Geometric invariant theory}, Springer-Verlag, 1965; 第 2 版
  与 J.~Fogarty 合著, 1982; 第 3 版与 J.~Fogarty 及 F.~Kirwan 合著, 1994.
\item[{[Mur65]}]\phantomsection\label{bib:sga3-V-Mur65} J.~P. Murre,
  \emph{Representation of unramified functors. Applications}
  (根据 A.~Grothendieck 未发表的结果), S\'em.
  Bourbaki, 第~9 卷, 第~294 篇 (1965), Soc. Math. France, 1995.
\item[{[Ray67a]}]\phantomsection\label{bib:sga3-V-Ray67a} M.~Raynaud,
  \emph{Passage au quotient par une relation d'\'equivalence plate},
  载于 \emph{Proc. Conf. Local Fields (Driebergen)}, 第~78--85 页
  (T.~A. Springer 编), Springer-Verlag, 1967.
\item[{[Ray67b]}]\phantomsection\label{bib:sga3-V-Ray67b} M.~Raynaud,
  \emph{Sur le passage au quotient par un groupo\"ide plat},
  C. R. Acad. Sci. Paris (S\'er. A) \textbf{265} (1967), 384--387.
\end{description}

% 翻译统计: 源行数 203; 输出行数 203 (含本统计行); 行内/行间数学片段 73; 翻译占位 0. 底本疑点另报主线程.
